%%%%%%%%%%%%%%%%%%%%%%%%%%%%%%%%%%%%%%%%%%%%%%%%%%%%%%%%%%%
% Matrix Class
%%%%%%%%%%%%%%%%%%%%%%%%%%%%%%%%%%%%%%%%%%%%%%%%%%%%%%%%%%%
%%%%%%%%%%%%%%%%%%%
% Matrix Creation %
%%%%%%%%%%%%%%%%%%%
% Create a matrix from values in shape of {{xx, xy, xz}{yx, yy, yz}}
\newcommand\newMat[2]{
    \xintFor* ##1 in #2 \do {%
        \xintifForFirst{%
            \newMatRow{matRow}{##1}
            \expandafter\edef\csname #1\endcsname{\xintCSVtoList{\matRow}}
        }{%
            \newMatRow{matRow}{##1}
            \appendToMat{#1}{\csname #1\endcsname}{\matRow}
        }
    }
}

% Create a matrix from row vectors
% TODO actually this takes matRows and not vecs
\newcommand\newMatFromRowVec[2]{
    \FPset\numParam{\expandafter\xintLength{#2}}%
    \FPsub\numParam\numParam{1}
    \FPdiv\numParam\numParam{2}
    \FPadd\numParam\numParam{1}
    % In case only one item/row was given the
    % matrix-datastructure must be adjusted
    \FPifeq{\numParam}{1}
        \expandafter\edef\csname #1\endcsname{\xintCSVtoList{#2}}%
    \else
        \expandafter\edef\csname #1\endcsname{\xintCSVtoList{#2}}%
    \fi
}

\newcommand\newEmptyMat[1]{
    \expandafter\edef\csname #1\endcsname{}%
}

% TODO adjust to variable length
\newcommand\newDiagMat[2]{
    \edef\tmpVector{\xintCSVtoList{#2}}%
    % \expandafter\edef\csname #1\endcsname{}%
    \newMatRow{rowA}{\xintNthElt{1}{\tmpVector}, 0, 0}
    \newMatRow{rowB}{0, \xintNthElt{2}{\tmpVector}, 0}
    \newMatRow{rowC}{0, 0, \xintNthElt{3}{\tmpVector}}
    \newMatFromRowVec{#1}{\rowA, \rowB, \rowC}
}

\newcommand\newMatRow[2]{
    \expandafter\edef\csname #1\endcsname{\xintCSVtoList{#2, 1}}%
}

\newcommand\appendToMat[3]{% {Basename}{input}
    \expandafter\edef\csname #1\endcsname{#2 \xintCSVtoList{#3}}%
}

\newcommand\vecToMatRow[2]{
    \expandafter\edef\csname #1\endcsname{#2 \xintCSVtoList{1}}%
}

% TODO BIG !!!
% THIS might not be needed anymore (sheet04 uses it)
% TODO automatically do this depending on the function
% or only have one representation (matrix) and adjust the functions to that
\newcommand\matColToVec[2]{
    \edef\numElems{#2}
    \edef\colLength{\xintSeq[+1]{+1}{\expandafter\xintLength\expandafter{\numElems}}}%
    \xintFor* ##1 in \colLength \do {%
        \xintifForFirst{%
            \newVec{testVec}{\xintNthElt{1}{\xintNthElt{1}{#2}}}
        }{%
            \appendToVec{testVec}{\testVec}{\xintNthElt{1}{\xintNthElt{##1}{#2}}}
        }
    }
    \copyVec[#1]{\testVec}
}

%%%%%%%%%%%%%%%%%%%%%%%%%%%
% Transformation Matrices %
%%%%%%%%%%%%%%%%%%%%%%%%%%%
% Create translation matrix from given vector
% #1 out, var name to save the result (matrix) to
% #2 a, translation vector
\newcommand\newTransMat[2]{%
    \edef\vecLength{\expandafter\xintLength\expandafter{#2 1}}
    \edef\innersequence{\xintSeq[+1]{+1}{\expandafter\xintLength\expandafter{#2}+1}}
    \xintFor* ##1 in \innersequence \do {%
        \xintifForFirst{%
            \xintFor* ##2 in \innersequence \do {%
                \newScalar{tmp}{0}
                \FPifeq{##2}{\vecLength}
                    \setScalar{tmp}{\xintNthElt{##1}{#2}}
                \else
                \fi
                \FPifeq{##1}{##2}
                    \setScalar{tmp}{1}
                \else
                \fi
                \xintifForFirst{%
                    \newVec{matRow}{\tmp}
                }{%
                    \appendToVec{matRow}{\matRow}{\tmp}
                }
            }
            \vecToMatRow{matRowB}{\matRow}
            \newMatFromRowVec{#1}{\matRowB}
        }{%
            \xintFor* ##2 in \innersequence \do {%
                \newScalar{tmp}{0}
                \FPifeq{##2}{\vecLength}
                    \setScalar{tmp}{\xintNthElt{##1}{#2}}
                \else
                \fi
                \FPifeq{##1}{##2}
                    \setScalar{tmp}{1}
                \else
                \fi
                \xintifForFirst{%
                    \newVec{matRow}{\tmp}
                }{%
                    \appendToVec{matRow}{\matRow}{\tmp}
                }
            }
            \vecToMatRow{matRowB}{\matRow}
            \appendToMat{#1}{\csname #1\endcsname}{\matRowB}
        }
    }
}

% Create scaling matrix from given vector
% #1 out, var name to save the result (matrix) to
% #2 a, scaling vector
\newcommand\newScaleMat[2]{%
    \edef\vecLength{\expandafter\xintLength\expandafter{#2 1}}
    \edef\innersequence{\xintSeq[+1]{+1}{\expandafter\xintLength\expandafter{#2}+1}}
    \xintFor* ##1 in \innersequence \do {%
        \xintifForFirst{%
            \xintFor* ##2 in \innersequence \do {%
                \FPiflt{##2}{\vecLength}
                    \setScalar{tmp}{\xintNthElt{##1}{#2}}
                \else
                    \setScalar{tmp}{1}
                \fi
                \FPifeq{##1}{##2}
                \else
                    \setScalar{tmp}{0}
                \fi
                \xintifForFirst{%
                    \newVec{matRow}{\tmp}
                }{%
                    \appendToVec{matRow}{\matRow}{\tmp}
                }
            }
            \vecToMatRow{matRowB}{\matRow}
            \newMatFromRowVec{#1}{\matRowB}
        }{%
            \xintFor* ##2 in \innersequence \do {%
                \FPiflt{##2}{\vecLength}
                    \setScalar{tmp}{\xintNthElt{##1}{#2}}
                \else
                    \setScalar{tmp}{1}
                \fi
                \FPifeq{##1}{##2}
                \else
                    \setScalar{tmp}{0}
                \fi
                \xintifForFirst{%
                    \newVec{matRow}{\tmp}
                }{%
                    \appendToVec{matRow}{\matRow}{\tmp}
                }
            }
            \vecToMatRow{matRowB}{\matRow}
            \appendToMat{#1}{\csname #1\endcsname}{\matRowB}
        }
    }
}

% Create rotation matrix from given scalar (2D)
% #1 out, var name to save the result (matrix) to
% #2 a, rotation value (scalar)
\newcommand\newRotMatTwoD[2][rotMatResult]{%
    \FPcos\tmpCos{#2}
    \FPsin\tmpSin{#2}
    \newMatRow{rowB}{\tmpSin, \tmpCos, 0}
    \FPmul\tmpSin\tmpSin{-1}
    \newMatRow{rowA}{\tmpCos, \tmpSin, 0}
    \newMatRow{rowC}{0, 0, 1}
    \newMatFromRowVec{#1}{\rowA, \rowB, \rowC}
}

% #1 Return Value Name
% #2 translation in x
% #3 translation in y
% \newcommand\newTransMatTwoD[3][transMatResult]{%
%     \newMatRow{rowA}{1, 0, #2}
%     \newMatRow{rowB}{0, 1, #3}
%     \newMatRow{rowC}{0, 0, 1}
%     \newMat{#1}{\rowA, \rowB, \rowC}
% }

% clone
% copy
% create

% fromMat2d
% fromMat4
% fromQuat
% fromQuat2
% fromRotation
% fromRotationTranslation
% fromRotationTranslationScale
% fromRotationTranslationScaleOrigin
% fromScaling
% fromTranslation
% fromXRotation
% fromYRotation
% fromZRotation

% frustum
% lookAt
% ortho
% orthoNO
% orthoZO
% perspective
% perspectiveFromFieldOfView
% perspectiveNO
% perspectiveZO

% normalFromMat4

% #1 Return Value Name
% #2 right vector
% #3 up vector
% #4 direction vector
% #5 camera position
% TODO
\newcommand\newLookAtMat[5][lookAtMatResult]{%
    \dotVec[scalarRC]{#2}{#5}
    \mulScalar[scalarRC]{\scalarRC}{-1}
    \copyVec[tmpVecA]{#2}
    \appendToVec{tmpVecA}{\tmpVecA}{\scalarRC}
    % \printVec[1]{\tmpVecA}
    %
    \dotVec[scalarUC]{#3}{#5}
    \mulScalar[scalarUC]{\scalarUC}{-1}
    \copyVec[tmpVecB]{#3}
    \appendToVec{tmpVecB}{\tmpVecB}{\scalarUC}
    %
    \dotVec[scalarDC]{#4}{#5}
    \copyVec[tmpVecC]{#4}
    \negateVec[tmpVecC]{\tmpVecC}
    \appendToVec{tmpVecC}{\tmpVecC}{\scalarDC}
    %
    \vecToMatRow{vecA}{\tmpVecA}
    \vecToMatRow{vecB}{\tmpVecB}
    \vecToMatRow{vecC}{\tmpVecC}
    \newMatRow{vecD}{0, 0, 0, 1}
    \newMatFromRowVec{#1}{\vecA, \vecB, \vecC, \vecD}
}

\newcommand\newProjMat[1]{%
    \newMatRow{vecA}{1, 0, 0, 0}
    \newMatRow{vecB}{0, 1, 0, 0}
    \newMatRow{vecC}{0, 0, 1, 0}
    \newMatRow{vecD}{0, 0, -1, 0}
    \newMatFromRowVec{#1}{\vecA, \vecB, \vecC, \vecD}
}

% TODO
\newcommand\stdProject[2]{%
    \newMatRow{vecA}{1, 0, 0, 0}
    \newMatRow{vecB}{0, 1, 0, 0}
    \newMatRow{vecC}{0, 0, 1, 0}
    \newMatRow{vecD}{0, 0, -1, 0}
    \newMat{matPstd}{\vecA, \vecB, \vecC, \vecD}
    \mulMatAndVec[#1]{\matPstd}{#2}
}

% #1 result var name
% #2 fovx
% #3 fovy
% #4 near Plane
% #5 far Plane
% TODO
\newcommand\newFrustumMat[5][frustMatResult]{%
    \FPdiv\tmp{#2}{2}
    \FPtan\frustumRight\tmp
    \FPmul\frustumRight{#4}\frustumRight
    \FPmul\frustumLeft\frustumRight{-1}
    \FPdiv\tmp{#3}{2}
    \FPtan\frustumTop\tmp
    \FPmul\frustumTop{#4}\frustumTop
    \FPmul\frustumBottom\frustumTop{-1}
    %
    \FPmul\cellAA{2}{#4} % AA
    \FPsub\tmp\frustumRight\frustumLeft
    \FPdiv\cellAA\cellAA\tmp
    \FPadd\cellAC\frustumRight\frustumLeft % AC
    \FPdiv\cellAC\cellAC\tmp
    \FPmul\cellBB{2}{#4} % BB
    \FPsub\tmp\frustumTop\frustumBottom
    \FPdiv\cellBB\cellBB\tmp
    \FPadd\cellBC\frustumTop\frustumBottom % BC
    \FPdiv\cellBC\cellBC\tmp
    \FPmul\cellCC{#5}{-1} % CC
    \FPsub\cellCC\cellCC{#4}
    \FPsub\tmp{#5}{#4}
    \FPdiv\cellCC\cellCC\tmp
    \FPmul\cellCD{#5}{#4} % CD
    \FPmul\cellCD{-2}\cellCD
    \FPdiv\cellCD\cellCD\tmp
    %
    \newMatRow{rowA}{\cellAA, 0, \cellAC, 0}
    \newMatRow{rowB}{0, \cellBB, \cellBC, 0}
    \newMatRow{rowC}{0, 0, \cellCC, \cellCD}
    \newMatRow{rowD}{0, 0, -1, 0}
    \newMatFromRowVec{#1}{\rowA, \rowB, \rowC, \rowD}
}

%%%%%%%%%%%%%%%%%%%
% Matrix Printing %
%%%%%%%%%%%%%%%%%%%
\newcommand\printMat[2][0]{%
    \ensuremath{%
        \begin{bmatrix}
            \xintFor* ##1 in #2 \do {%
                \xintFor* ##2 in {##1} \do {%
                    \xintifForLast{%
                        \\
                    }{%
                        \xintifForFirst{}{&}
                        \FPifeq{#1}{0}
                            \roundScalar{##2}
                            \scalarRoundResult
                        \else
                            \numToFractionA{##2}
                        \fi
                    }
                }
            }
        \end{bmatrix}
    }
}

% TODO variable length
\newcommand\printMatDiag[2][0]{%
    \ensuremath{%
        \text{diag}\left(
        \xintNthElt{1}{\xintNthElt{1}{#2}},
        \xintNthElt{2}{\xintNthElt{2}{#2}},
        \xintNthElt{3}{\xintNthElt{3}{#2}}
        \right)
    }
}

\newcommand\getMatHeight[2][matHeightResult]{%
    \newScalar{#1}{\expandafter\xintLength\expandafter{#2}}%
}

\newcommand\getMatWidth[2][matWidthResult]{%
    \edef\rowElements{\xintNthElt{1}{#2}}%
    \newScalar{#1}{\expandafter\xintLength\expandafter{\rowElements}}%
    \subScalar[#1]{\csname #1\endcsname}{1}%
    %
    % \newScalar{\matWidthResult}{4}%
    % \FPset{\matWidthResultA}{4}
    % \FPsub{\matWidthResultC}{\matWidthResultA}{1}
    % \FPset{\matWidthResult}{\matWidthResultC}
    %
    % \edef\matWidthResult{\matWidthResultC}%
    % \subScalar[#1]{\csname #1\endcsname}{1}%
    %
    % % \expandafter\edef\csname #1\endcsname{\expandafter\xintLength\expandafter{\rowElements}}%
    % \newScalar{tmp}{\expandafter\xintLength\expandafter{\rowElements}}%
    % \FPsub\tmp{\tmp}{1}
    % % \newScalar{#1}{\tmp}
    % \edef\matWidthResult{\tmp}%
    % % \subScalar[#1]{\csname #1\endcsname}{1}%
    % % \newScalar{#1}{\aaa}%
}

%%%%%%%%%%%%%%%%%%%%%%%
% Matrix Modification %
%%%%%%%%%%%%%%%%%%%%%%%

% translate
% transpose
% rotate
% rotateX
% rotateY
% rotateZ
% scale
% set
% projection
% targetTo

%%%%%%%%%%%%%%%
% Matrix Math %
%%%%%%%%%%%%%%%
% Multiply matrix with a scalar
% #1 out, var name to save the result (matrix) to
% #2 a, matrix value
% #3 b, scalar value
\newcommand\mulScalarMat[3][mulMatScalarResult]{%
    \edef\elemOne{\xintNthElt{1}{#2}}
    \edef\innersequence{\xintSeq[+1]{+1}{\expandafter\xintLength\expandafter{\elemOne}-1}}%
    %
    \xintFor* ##1 in #2 \do {%
        \xintifForFirst{%
            \xintFor* ##2 in \innersequence \do {%
                \FPset\elem{\xintNthElt{##2}{##1}}
                \FPmul\elem\elem{#3}
                \xintifForFirst{%
                    \newVec{rowVec}{\elem}
                }{%
                    \appendToVec{rowVec}{\rowVec}{\elem}
                }
            }
            \vecToMatRow{matRow}{\rowVec}
            \newMatFromRowVec{#1}{\matRow}
        }{%
            \xintFor* ##2 in \innersequence \do {%
                \FPset\elem{\xintNthElt{##2}{##1}}
                \FPmul\elem\elem{#3}
                \xintifForFirst{%
                    \newVec{rowVec}{\elem}
                }{%
                    \appendToVec{rowVec}{\rowVec}{\elem}
                }
            }
            \vecToMatRow{matRow}{\rowVec}
            \appendToMat{#1}{\csname #1\endcsname}{\matRow}
        }
    }
}

\newcommand\mulMat[3][mulMatResult]{%
    \edef\elemsRow{\xintNthElt{1}{#3}}
    \edef\widthsequence{\xintSeq[+1]{+1}{\expandafter\xintLength\expandafter{\elemsRow}}}%
    %
    \edef\elemsMat{#2}
    \edef\heightsequence{\xintSeq[+1]{+1}{\expandafter\xintLength\expandafter{\elemsMat}}}%
    %
    \edef\elemsMat{#3}
    \edef\heightsequenceB{\xintSeq[+1]{+1}{\expandafter\xintLength\expandafter{\elemsMat}}}%
    %
    \edef\elemsRowTmp{\xintNthElt{1}{#2}}
    \FPset\widthMatA{\expandafter\xintLength\expandafter{\elemsRowTmp}}%
    \FPsub\widthMatA\widthMatA{1}
    \FPset\heightMatB{\expandafter\xintLength\expandafter{\elemsMat}}%
    % TODO check this
    % \FPifeq{\widthMatA}{\heightMatB}
    % \else
    %     \PackageError{glmatrix package}{Dimension mismatch for matrix, could not multiply}{width (\widthMatA) of matrix #2 and height (\heightMatB) of matrix #3 do not match}
    % \fi
    %
    \newEmptyMat{#1}
    \xintFor* ##2 in \heightsequence \do {%
        \xintFor* ##3 in \widthsequence \do {%
            \FPset\value{0}
            \xintifForFirst{%
                \xintFor* ##4 in \heightsequenceB \do {%
                    \FPmul\tmptmp{%
                    \xintNthElt{##4}{\xintNthElt{##2}{#2}}%
                    }{%
                    \xintNthElt{##3}{\xintNthElt{##4}{#3}}%
                    }
                    \FPadd\value\value\tmptmp
                }
                % }
                \newVec{rowVec}{\value}
            }{%
                \xintifForLast{
                }{
                    \xintFor* ##4 in \heightsequenceB \do {%
                        \FPmul\tmptmp{%
                        \xintNthElt{##4}{\xintNthElt{##2}{#2}}%
                        }{%
                        \xintNthElt{##3}{\xintNthElt{##4}{#3}}%
                        }
                        \FPadd\value\value\tmptmp
                    }
                    \appendToVec{rowVec}{\rowVec}{\value}
                }
            }
        }
        \vecToMatRow{matRow}{\rowVec}
        \appendToMat{#1}{\csname #1\endcsname}{\matRow}
    }
    % Test if width and height is correct
    \edef\expectedheight{\expandafter\xintLength\expandafter{#2}}%
    \edef\elemsRow{\xintNthElt{1}{#3}}
    \edef\expectedwidth{\expandafter\xintLength\expandafter{\elemsRow}}%
    %
    \edef\elemsRow{\xintNthElt{1}{\csname #1\endcsname}}
    \edef\actualwidth{\expandafter\xintLength\expandafter{\elemsRow}}%
    \edef\elemsMat{\csname #1\endcsname}
    \edef\actualheight{\expandafter\xintLength\expandafter{\elemsMat}}%
    %
    \FPifeq{\expectedwidth}{\actualwidth}
    \else
        \PackageError{glmatrix package}{Matrix width does not match expectation}{width seems to be incorrect}
    \fi
    \FPifeq{\expectedheight}{\actualheight}
    \else
        \PackageError{glmatrix package}{Matrix height does not match expectation}{height seems to be incorrect}
    \fi
}

% TODO BIG !!!
\newcommand\mulMatAndVec[3][mulMatResult]{%
    \newMatRow{vecA}{\xintNthElt{1}{#3}}
    \newMatRow{vecB}{\xintNthElt{2}{#3}}
    \newMatRow{vecC}{\xintNthElt{3}{#3}}
    \newMatFromRowVec{matB}{\vecA, \vecB, \vecC}
    \mulMat[#1]{#2}{\matB}
}

% add
% mul multiply
% multiplyScalar
% multiplyScalarAndAdd
% sub subtract

% Sum list of matrizes
% TODO
\newcommand\addMat[2][sumMatResult]{%
    \def\matlist{\xintCSVtoList{#2}}
    \edef\elemsRow{\xintNthElt{1}{\xintNthElt{1}{\matlist}}}
    \edef\widthsequence{\xintSeq[+1]{+1}{\expandafter\xintLength\expandafter{\elemsRow}}}%
    \edef\elemsMat{\xintNthElt{1}{\matlist}}
    \edef\heightsequence{\xintSeq[+1]{+1}{\expandafter\xintLength\expandafter{\elemsMat}}}%
    %
    \newEmptyMat{#1}
    \xintFor* ##2 in \heightsequence \do {%
        \xintFor* ##3 in \widthsequence \do {%
            \FPset\value{0}
            \xintifForFirst{%
                \xintFor* ##1 in \matlist \do {%
                    \FPadd\value\value{\xintNthElt{##3}{\xintNthElt{##2}{##1}}}
                }
                \newVec{rowVec}{\value}
            }{%
                \xintifForLast{
                }{
                    \xintFor* ##1 in \matlist \do {%
                        \FPadd\value\value{\xintNthElt{##3}{\xintNthElt{##2}{##1}}}
                    }
                    \appendToVec{rowVec}{\rowVec}{\value}
                }
            }
        }
        \vecToMatRow{matRow}{\rowVec}
        \appendToMat{#1}{\csname #1\endcsname}{\matRow}
    }
}

%%%%%%%%%%%%%%%%%%%%%
% Matrix Comparison %
%%%%%%%%%%%%%%%%%%%%%
% equals
% getRotation
% getScaling
% getTranslation

%%%%%%%%%%%%%%%
% Matrix Misc %
%%%%%%%%%%%%%%%
% adjoint
% decompose
% determinant
% frob
% identity
% invert

% LDU
% str
